\section{Introduction}
\label{sec:introduction}

Quantum imaging harnesses the non-classical correlations of entangled or correlated photon pairs to achieve imaging modalities that are not accessible with classical light sources. Since the seminal demonstration of imaging via two-photon entanglement by \citet{pittman1995}, the field has expanded to encompass ghost imaging, quantum illumination, sub-shot-noise microscopy, and entanglement-assisted compressed sensing. These techniques promise advantages in spatial resolution, noise robustness, and the ability to image with photons that never interact with the object \citep{erkmen2010}.

Despite these capabilities, a key challenge remains: \emph{image reconstruction}. Quantum imaging measurements are inherently sparse, limited by detector efficiencies, coincidence counting rates, and photon budgets, and the resulting inverse problems are severely ill-posed. Traditional reconstruction approaches rely on correlation computation or iterative optimisation, which demand large numbers of measurements and significant computation time.

Deep learning has proved effective for solving inverse problems in computational imaging \citep{barbastathis2019, lecun2015} and is now being applied to quantum imaging reconstruction. Neural networks can learn strong image priors from data, enabling high-quality reconstruction from far fewer measurements than classical algorithms require. Work at the intersection of quantum optics and deep learning now spans multiple imaging modalities and network architectures.

\subsection{Scope and Contributions}

This survey provides a systematic review of deep learning methods applied to quantum imaging reconstruction. We focus on three principal modalities:

\begin{enumerate}[label=(\roman*)]
    \item \textbf{Quantum ghost imaging}, where deep learning dramatically reduces measurement requirements for image reconstruction from spatially correlated photon fields.
    \item \textbf{Compressed sensing with quantum correlations}, where neural networks enhance reconstruction from sub-Nyquist quantum measurements.
    \item \textbf{Sub-shot-noise imaging}, where deep learning denoising exploits quantum noise statistics to push sensitivity below the classical shot-noise limit.
\end{enumerate}

The contributions of this survey are:
\begin{itemize}
    \item A unified taxonomic framework classifying methods by imaging modality, DL architecture, and reconstruction task.
    \item A detailed comparative analysis of deep learning approaches for ghost imaging reconstruction, including quantitative performance comparisons.
    \item Identification of critical gaps in benchmarking, reproducibility, and standardisation.
    \item A discussion of hybrid classical--quantum architectures and future research directions.
\end{itemize}

\noindent Figure~\ref{fig:taxonomy} illustrates this classification framework.

\begin{figure*}[t]
\centering
\begin{tikzpicture}[
    node distance=0.4cm and 0.6cm,
    rootnode/.style={rectangle, rounded corners, draw=black!80, fill=black!8, font=\bfseries\small, minimum height=0.7cm, minimum width=2.8cm, align=center},
    axisnode/.style={rectangle, rounded corners, draw=black!70, fill=blue!8, font=\small\bfseries, minimum height=0.6cm, minimum width=2.2cm, align=center},
    leafnode/.style={rectangle, rounded corners, draw=black!50, fill=white, font=\footnotesize, minimum height=0.5cm, minimum width=1.8cm, align=center},
    arr/.style={-{Stealth[length=2.5mm]}, thick, black!60}
]

% Root
\node[rootnode] (root) {Survey Taxonomy};

% Three axes
\node[axisnode, below left=0.7cm and 3.2cm of root] (mod) {Imaging Modality};
\node[axisnode, below=0.7cm of root] (arch) {DL Architecture};
\node[axisnode, below right=0.7cm and 3.2cm of root] (task) {Reconstruction Task};

\draw[arr] (root) -- (mod);
\draw[arr] (root) -- (arch);
\draw[arr] (root) -- (task);

% Modality leaves
\node[leafnode, below=0.25cm of mod, xshift=-1.1cm] (m1) {Ghost Imaging};
\node[leafnode, right=0.15cm of m1] (m2) {Compressed\\[-1pt]Sensing};
\node[leafnode, below=0.25cm of mod, xshift=1.1cm] (m3) {Sub-Shot-\\[-1pt]Noise};
\node[leafnode, below=0.35cm of m2] (m4) {Hybrid};

\draw[arr, thin] (mod) -- (m1);
\draw[arr, thin] (mod) -- (m2);
\draw[arr, thin] (mod) -- (m3);
\draw[arr, thin] (mod) -- (m4);

% Architecture leaves
\node[leafnode, below=0.25cm of arch, xshift=-1.1cm] (a1) {CNN / U-Net};
\node[leafnode, right=0.15cm of a1] (a2) {GAN};
\node[leafnode, below=0.25cm of arch, xshift=1.1cm] (a3) {Deep\\[-1pt]Unfolding};
\node[leafnode, below=0.35cm of a2] (a4) {Quantum-\\[-1pt]Inspired};

\draw[arr, thin] (arch) -- (a1);
\draw[arr, thin] (arch) -- (a2);
\draw[arr, thin] (arch) -- (a3);
\draw[arr, thin] (arch) -- (a4);

% Task leaves
\node[leafnode, below=0.25cm of task, xshift=-1.1cm] (t1) {Compressed\\[-1pt]Reconstruction};
\node[leafnode, right=0.15cm of t1] (t2) {Denoising};
\node[leafnode, below=0.35cm of t1, xshift=1.1cm] (t3) {End-to-End};

\draw[arr, thin] (task) -- (t1);
\draw[arr, thin] (task) -- (t2);
\draw[arr, thin] (task) -- (t3);

\end{tikzpicture}
\caption{Taxonomic classification framework used in this survey. Methods are categorised along three axes: the quantum imaging modality, the deep learning architecture, and the reconstruction task.}
\label{fig:taxonomy}
\end{figure*}

\subsection{Paper Organisation}

Section~\ref{sec:background} reviews the necessary background in quantum imaging and deep learning for inverse problems. Section~\ref{sec:ghost_imaging} surveys deep learning methods for ghost imaging in detail. Section~\ref{sec:compressed_sensing} covers compressed sensing with quantum correlations. Section~\ref{sec:sub_shot_noise} addresses sub-shot-noise imaging. Section~\ref{sec:hybrid} examines hybrid classical--quantum architectures. Section~\ref{sec:discussion} discusses open problems, and Section~\ref{sec:conclusion} concludes.
